\chapter{Literature Review}
\label{ch:literature}

\todo{this chapter needs the shortlisted papers read and summarised; the
candidate list is in docs/04-reference/paper-shortlist.md}

\section{Manufacturing Resource Planning}

\todo{MRP to MRP II to ERP: what each layer added, and where central planning
becomes a weakness}

\section{Multi-Agent Systems in Manufacturing}

\todo{survey of agent-based manufacturing control; identify which papers
distribute scheduling and which only distribute monitoring}

\section{Negotiation Protocols}

\subsection{The Contract Net Protocol}

The Contract Net Protocol (CNP) assigns tasks by announcement and
bidding~\cite{smith1980contractnet}. A manager announces a task, contractors
submit bids, the manager awards the task to a chosen bid, and the winner
reports the result. The FIPA specification formalises this as a sequence of
call-for-proposal, proposal or refusal, acceptance or rejection, and finally
result or failure~\cite{fipa2002contractnet}. This project uses a simplified
CNP in which an order agent is the manager and resource agents are the
contractors.

\todo{state precisely which parts of FIPA compliance this project does and
does not claim --- examiners ask this}

\subsection{The Agent2Agent Protocol}

\todo{what A2A provides, how it differs from FIPA-ACL over XMPP, and why it was
chosen; see docs/03-planning/technology-decision.md}

\section{Agent Frameworks and Runtimes}

\todo{compare the frameworks evaluated -- Mesa, SPADE, PADE, and the
TypeScript/A2A approach adopted; summarise the comparison table from the
technology decision document}

\section{Language Models in Planning Systems}

\todo{why language models are unsuitable for the arithmetic and feasibility
parts of planning, and what they are appropriate for; this justifies the
project's rule that the model never computes or mutates state}

\section{Summary and Research Gap}

\todo{one paragraph: what existing work covers, what it leaves open, and how
this project sits in that space}
